\documentclass{article}
\usepackage[utf8]{inputenc}
\usepackage[T1]{fontenc}
\usepackage{geometry}
\geometry{
    a4paper,
    left=2.5cm,
    right=2.5cm,
    top=2.5cm,
    bottom=2.5cm
}
\usepackage{graphicx}
\usepackage{float}
\usepackage{hyperref}
\usepackage[english]{babel}
\usepackage{amsmath}
\usepackage{amssymb}
\usepackage{array}
\usepackage{booktabs}
\usepackage{caption}
\usepackage{subcaption}
\usepackage{enumitem}
\usepackage{fancyhdr}
\usepackage{titlesec}
\usepackage{url}
\usepackage{textcomp}
\usepackage{xcolor}
\usepackage{listings}
\usepackage{multirow}
\usepackage{longtable}
\usepackage{tabularx}
\usepackage{times}

% Set up page style
\pagestyle{fancy}
\fancyhf{}
\rfoot{\thepage}
\renewcommand{\headrulewidth}{0pt}

% Title formatting
\titleformat{\chapter}[display]{\normalfont\bfseries}{\chaptertitle\thechapter}{0pt}{\Huge}
\titlespacing*{\chapter}{0pt}{50pt}{40pt}

% Document info
\title{Template Document}
\author{Academic Support Agent}
\date{\today}

\begin{document}
\maketitle

%\layout

% DEBUT DE L'ARTICLE
%
\section{Introduction}

Pour produire un article en \LaTeX\ nous recommandons d'utiliser la
classe \textit{rnti.cls} version~2. Voir  galement le guide
d'utilisation de la classe \citep{ritschard:2005rnti}. Merci de
v rifier qu'il ne subsiste pas d'\emph{overfull box} (d passement
des marges).

Pour les textes de 2 pages ou moins (posters), ne pas mettre de
r sum  au d but. On donnera cependant un ``summary'' (en anglais)
d'au plus 40 mots en fin de texte. Utiliser pour cela l'option
\textit{noresume} qui supprime le r sum  mais maintient le summary.

\section{Consignes pour le texte}


Le texte est en Times 10 points (faites appel   l'extension
\textit{times.sty} si vous n'utilisez pas la classe \textit{rnti}).
Le texte de chaque page y compris les ent tes et pieds de pages doit
appartenir   une zone de 19.3 cm par 13.2 cm centr e horizontalement
et verticalement sur la page.

Vous trouverez ci-dessous une liste de petites erreurs qui reviennent souvent et qui nuisent   la qualit  des actes~:
\begin{itemize}
\item mettre un espace ins cable (avec le caract re \verb+~+) avant une commande \verb+\ref{}+ ou   \verb+\cite{}+,
\item  Utiliser les caract res  \verb+``+ et \verb+''+ pour produire des guillemets  corrects,
\item pas de \verb+\vspace{}+ n gatif pour modifier l'espacement entre les paragraphes, les 
parties de l'article, la l gende d'un tableau ou d'une figure, 
\item ne pas changer la taille de la bibliographie,
\item utiliser plut t une taille relative pour la taille des images en respectant 
le ratio hauteur/largeur,
\item pas d'overfull hbox,
\item  viter d' crire ... mais utiliser plut t les commandes \verb+\dots+, \verb+\ldots+  ou \verb+\cdots+ selon l'effet recherch ,
\item si les auteurs ont tous la m me adresse (une seule adresse sur la premi re page de l'article), ne pas mettre d' toile apres les noms des auteurs et avant l'adresse.
\end{itemize}

\subsection{Section de niveau 2}

Exemple de section de niveau 2.


\subsubsection{Section de niveau 3}

Exemple de section de niveau 3. Il vaut mieux  viter, dans la mesure
du possible, d'utiliser ce niveau de section surtout pour les
articles courts. Utiliser de pr f rence
\verb+\paragraph{+\textit{Titre du paragraphe.}\verb+}+ qui produit

\paragraph{Titre du paragraphe.} Ceci est un exemple de section type
paragraphe. On note l'absence de num rotation.

\subsection{Liste d' l ments}

Merci de respecter les styles suivants pour les listes.

Exemple de liste sans num rotation :
\begin{itemize}
    \item premier item ;
    \item deuxi me item ;
    \begin{itemize}
      \item premier sous-item,
      \item second sous-item,
      \item dernier sous-item.
  \end{itemize}
  \item dernier item.
\end{itemize}

Autre liste possible avec num rotation :
\begin{enumerate}
    \item premier item ;
    \item deuxi me item ;
    \begin{enumerate}
      \item premier sous-item,
      \item second sous-item,
      \item dernier sous-item.
    \end{enumerate}
  \item dernier item.
\end{enumerate}

\subsection{Tables, Figures}


Les tableaux et figures doivent  tre centr s horizontalement. Les
l gendes sont plac es sous ces m mes  l ments et sont centr es si
elle font moins d'une ligne et justifi es sinon. Voir le
tableau~\ref{tab_exemple} et la figure~\ref{fig_exemple} pour des
exemples. Les l gendes doivent se terminer avec un point.



\begin{table}[ht]
 \begin{center}
   \tabcolsep = 2\tabcolsep
   \begin{tabular}{lcc}
   \hline\hline
                & mod le 1 & mod le 2 \\
   \hline
   indicateur 1 & 1        & 2        \\
   indicateur 2 & 2        & 1        \\
   indicateur 3 & 1        & 1        \\
   \hline
   \end{tabular}
\caption{Exemple de tableau avec sa l gende. Merci de terminer la l gende par un point.} \label{tab_exemple}
 \end{center}
\end{table}

Pour les tableaux,  viter de multiplier les lignes s paratrices de
colonnes et de lignes, cela alourdi inutilement la page.


\begin{figure}[t]
\begin{center}
 \includegraphics[width=6cm]{fg_complex}
 \caption{Exemple de figure avec sa l gende. La l gende est centr e,
sauf si, comme ici, elle fait plus d'une ligne, auquel cas elle est
justifi e. Merci de terminer la l gende par un point.} \label{fig_exemple}
\end{center}
\end{figure}

N'oubliez pas, le cas  ch ant, d'indiquer la signification des axes
sur vos graphiques.


\subsection{Notes de bas de pages}

Le texte de la note appara t en bas de la page.\footnote{Exemple de
note de bas de page.}

\subsection{Noms des auteurs et titre court dans l'ent te}

Pour les noms des auteurs utilis s dans l'ent te de l'article, merci
de faire figurer un ou deux noms (avec initiales des pr noms) pour
des articles   un ou deux auteurs (``Initiale. NomPremierAuteur'',
ou ``Initiale1. NomPremierAuteur et Initiale2. NomDeuxi meAuteur''),
et pour les articles   trois auteurs ou plus, utiliser une notation
du type ``Initiale. NomPremierAuteur et al.''. Seule la premi re
lettre des noms doit  tre en capitale, le reste est en minuscules.
De m me, le titre court doit  tre en minuscules.



\subsection{Pas de pagination, page A4 et article final au format PDF}

Ne pas indiquer la num rotation des pages de votre article car
celle-ci sera ajout e lors de la fabrication du num ro.

Dans la version soumise pour  valuation, il est souhaitable
cependant de num roter les pages. Ceci facilite les commentaires des
relecteurs.

L'article final doit  tre imp rativement remis au format PDF avec
une dimension de page A4 (la taille letter n'est pas admise).
L'option \texttt{a4paper} produit du format A4 seulement si vous
compilez avec pdf\LaTeX{}. Si vous produisez un .dvi ou .ps que vous
covertissez ensuite en pdf avec par exemple dvi2pdf ou ps2pdf, il
vous appartient de configurer correctement votre convertisseur ou de
lui fournir l'option appropri e (-pagesize=A4).


\section{Consignes pour les r f rences}

Les r f rences sont donn es en fin d'article avant le
\guilo{}Summary\guilf{}. Elles doivent  tre list es par ordre
alphab tique.  Nous vous recommandons l'utilisation de
\textsc{Bib}\TeX\ avec le style bibliographique \textit{rnti.bst}
d crit dans \citet{ritschard:2005rnti}. Si vous n'utilisez pas ce
style, merci de suivre les exemples donn s   la fin de cet exemple
d'article.

Dans le corps du texte, on utilise \citet{thSau00}, \citet{HolWil90}
pour faire r f rence   un article avec un ou deux auteurs et
\citet{lioni01} lorsque trois auteurs ou plus sont pr sents.  Selon
les cas, on pourra aussi utiliser \citep{thSau00} ou
\citep[voir][chapitre 4]{brei:frie:olsh:ston:1984}, ou encore
\citet{quin:1986ID3,quin:1993} dans le cas de citations multiples
d'un m me auteur. Merci de suivre ces exemples.







\section*{Annexe}

Voici un exemple d'annexe. S'il y a plus de deux annexes, merci de
les num roter (Annexe~1, Annexe~2, etc).

\Fr
\end{document}